\documentclass[10pt,a4paper]{amsart}
\usepackage[margin=19mm]{geometry}
\usepackage{amsmath,amssymb,mathtools}
\usepackage[scr=rsfs,cal=euler]{mathalfa}
\usepackage{xfrac}
\usepackage[unicode,hidelinks,bookmarks=false]{hyperref}
\newcommand{\ie}{i.e.\ }
\newcommand{\cf}{cf.\ }
\newcommand{\resp}{respectively\ }
\newcommand{\Subsection}{\subsection}
\providecommand{\nobreakdash}{\nobreak\hbox{-}}
\setlength{\emergencystretch}{2em}
\multlinegap0pt
\usepackage{bm}
\usepackage{bbm}
\usepackage{dsfont}
\usepackage{xspace}
\usepackage{cancel}
\usepackage{mathtools}
\usepackage{amsthm}
\makeatletter
\@ifundefined{theorem}{\newtheorem{theorem}{Theorem}}{}
\@ifundefined{proposition}{\newtheorem{proposition}[theorem]{Proposition}}{}
\makeatother
\title[Transversal Difference Numbers in Finite Abelian Quotients]{Transversal Difference Numbers in Finite Abelian Quotients}
\author{Mugurel Barcau, Vicen\c{t}iu Pa\c{s}ol, George C. \c{T}urca\c{s}}
\date{Source: arXiv:2606.27961v1 (CC BY 4.0)}
\begin{document}
\maketitle

\noindent\emph{Source and licence.} Transversal Difference Numbers in Finite Abelian Quotients. Version arXiv:2606.27961v1 (CC BY 4.0), distributed under the Creative Commons Attribution 4.0 International licence, as recorded in the arXiv licence metadata for this exact version and shown on the abstract page https://arxiv.org/abs/2606.27961v1. Full rights notice: licenses/arxiv-2606.27961-CC-BY-4.0.txt.\par\smallskip

\noindent\textit{Editorial scope.} Counted unit: the Proposition whose source label is \texttt{prop:primitive-quotient-reduction} in the article (source0.tex lines 570--611, source label \texttt{prop:primitive-quotient-reduction}), delivered together with the article's own complete original proof of it (source0.tex lines 613--726, one \texttt{proof} environment of 114 lines). No mathematics is written, repaired or re-derived: statement and proof are the article's own text line for line. Required context retained uncounted: none. Every definition, display and label of those lines is retained unchanged. Omitted from this package and not redistributed: 1--569, 612--612, 727--2229. Nothing inside a retained excerpt is omitted or altered: every retained body is delivered exactly as the article's lines are written, so no omission receipt of any kind accompanies this package. Citations: the retained excerpts carry no citation command at all, so no bibliography entry of the article is transcribed and no reference is redistributed. Reference adaptation: none was needed and none was made; every reference inside the delivered unit resolves inside it, so no unresolved reference or citation is produced. Printed slips kept literal: no printed word, symbol, hypothesis or number of the counted statement or of its proof was altered. Layout and class: the article's own class is \texttt{amsart} with packages including amsrefs, amsmath, amsfonts, amssymb, amsthm, mathrsfs, xcolor, hyperref, comment, graphicx, caption, subcaption; the wrapper is the lane's pinned \texttt{amsart} editorial wrapper at 10pt on A4, which restates the theorem-like environments the retained text needs and adds the article's own macro definitions that the retained text needs (none), with extra wrapper packages (bm, bbm, dsfont, xspace, cancel, mathtools). The delivered numbering is the unit's own. Assets: no retained span contains a figure environment, a graphics-inclusion command, a PDF-page inclusion, a table or a TikZ picture; no image file is copied into this package, the package declares no asset dependency and no image file is redistributed with it. Licence: version arXiv:2606.27961v1 of this article is distributed under the Creative Commons Attribution 4.0 International licence, as recorded in the arXiv licence metadata for this exact version and shown on the abstract page https://arxiv.org/abs/2606.27961v1. A full rights notice is delivered at licenses/arxiv-2606.27961-CC-BY-4.0.txt. Characters: every retained excerpt is pure ASCII, so the delivered engine, XeLaTeX with its default Latin Modern text font, reports no missing character.
\begin{proposition}
\label{prop:primitive-quotient-reduction}
Let $H\leq G$, and let $T$ be a normalized transversal for $G/H$.  Then:
\begin{enumerate}
  \item $\Sigma(T)$ is a subgroup of $G/H$;

  \item the map
  \[
    \Sigma(T)\longrightarrow G,
    \qquad a\longmapsto d_T(a),
  \]
  is an injective group homomorphism, its image $\widetilde\Sigma(T)$ is a
  subgroup of $G$, and
  \[
    \widetilde\Sigma(T)\cap H=\{0\},
    \qquad
    |\widetilde\Sigma(T)|=|\Sigma(T)|;
  \]

  \item one has
  \[
    T+\widetilde\Sigma(T)=T,
    \qquad
    D(T)+\widetilde\Sigma(T)=D(T),
  \]
  so $D(T)$ is a union of cosets modulo $\widetilde\Sigma(T)$;

  \item if
  \[
    \overline G=G/\widetilde\Sigma(T),
    \qquad
    \overline H=(H+\widetilde\Sigma(T))/\widetilde\Sigma(T),
  \]
  and $\overline T$ denotes the image of $T$ in $\overline G$, then
  $\overline T$ is a transversal for $\overline G/\overline H$ and
  \[
    |D(T)|=|\widetilde\Sigma(T)|\,|D(\overline T)|;
  \]

  \item the transversal $\overline T$ has no nonzero singleton directions.
\end{enumerate}
\end{proposition}

\begin{proof}
Put $Q=G/H$ and write $s=s_T$ for the normalized section attached to $T$.
The zero direction is singleton, since $\mathcal D_0(T)=\{0\}$.  Let
$a,b\in\Sigma(T)$.  Then, for every $x\in Q$,
\[
  s(x+a)-s(x)=d_T(a),
  \qquad
  s(x+b)-s(x)=d_T(b).
\]
Hence
\[
\begin{aligned}
  s(x+a+b)-s(x)
  &=\bigl(s(x+a+b)-s(x+a)\bigr)+\bigl(s(x+a)-s(x)\bigr) \\
  &=d_T(b)+d_T(a).
\end{aligned}
\]
Thus $a+b$ is a singleton direction and
\[
  d_T(a+b)=d_T(a)+d_T(b).
\]
Since $Q$ is finite, closure under addition and the presence of $0$ imply that
$\Sigma(T)$ is a subgroup.  The displayed identity also shows that
$a\mapsto d_T(a)$ is a group homomorphism.  This proves (1) and the homomorphism
assertion in (2).

For $a\in\Sigma(T)$ we have
\[
  \pi(d_T(a))=a,
\]
because $d_T(a)=s(x+a)-s(x)$ for every $x\in Q$.  Hence the homomorphism
$a\mapsto d_T(a)$ is injective, and its image $\widetilde\Sigma(T)$ is a
subgroup of $G$ with $|\widetilde\Sigma(T)|=|\Sigma(T)|$.  If
$d_T(a)\in H$, then $a=\pi(d_T(a))=0$, hence $d_T(a)=0$.  Therefore
\[
  \widetilde\Sigma(T)\cap H=\{0\}.
\]
This proves (2).

If $\sigma=d_T(a)\in\widetilde\Sigma(T)$, then
\[
  s(x+a)=s(x)+\sigma
\]
for every $x\in Q$.  As $x$ varies over $Q$, so does $x+a$, and therefore
translation by $\sigma$ permutes the elements of $T$.  Thus
\[
  T+\widetilde\Sigma(T)=T.
\]
Subtracting gives
\[
  D(T)+\widetilde\Sigma(T)=D(T),
\]
so $D(T)$ is a union of cosets modulo $\widetilde\Sigma(T)$.

Let $q_{\natural}:G\to\overline G$ be the quotient map.  We prove (4).  The
image $\overline T=q_{\natural}(T)$ plainly meets every
$\overline H$-coset.  For uniqueness, suppose that $q_{\natural}(t_1)$ and
$q_{\natural}(t_2)$ lie in the same $\overline H$-coset.  Then
\[
  t_1-t_2\in H+\widetilde\Sigma(T).
\]
Write $t_1-t_2=h+\sigma$, with $h\in H$ and
$\sigma\in\widetilde\Sigma(T)$.  Since $T+\widetilde\Sigma(T)=T$, the element
$t_2+\sigma$ belongs to $T$.  But
\[
  t_1-(t_2+\sigma)=h\in H.
\]
Since $T$ is a transversal for $G/H$, this forces $t_1=t_2+\sigma$, hence
$q_{\natural}(t_1)=q_{\natural}(t_2)$.  Thus $\overline T$ is a transversal.

Moreover,
\[
  q_{\natural}(D(T))=q_{\natural}(T-T)=\overline T-\overline T=D(\overline T).
\]
Since $D(T)$ is a union of $\widetilde\Sigma(T)$-cosets, each element of
$D(\overline T)$ has exactly $|\widetilde\Sigma(T)|$ preimages in $D(T)$.  Hence
\[
  |D(T)|=|\widetilde\Sigma(T)|\,|D(\overline T)|.
\]

It remains to show that $\overline T$ is primitive.  Under the natural
identification
\[
  \overline G/\overline H
  \simeq
  G/(H+\widetilde\Sigma(T))
  \simeq
  (G/H)/\Sigma(T),
\]
suppose that a nonzero class $\overline a$ has singleton fibre for
$\overline T$, and choose a representative $a\in Q$.  Then, for some $d\in G$,
\[
  \mathcal D_a(T)\subseteq d+\widetilde\Sigma(T).
\]
Indeed, the quotient fibre in direction \(\overline a\) is the image of the
fibres \(\mathcal D_{a+k}(T)\), \(k\in\Sigma(T)\); if that image is a singleton,
then each such fibre lies in one \(\widetilde\Sigma(T)\)-coset.
For every $k\in\Sigma(T)$, the fibre in direction $a+k$ satisfies
\[
\begin{aligned}
  \mathcal D_{a+k}(T)
  &=\{s(x+a+k)-s(x):x\in Q\} \\
  &=\{s(x+a)+d_T(k)-s(x):x\in Q\} \\
  &=\mathcal D_a(T)+d_T(k).
\end{aligned}
\]
Thus the union of the fibres above $a+\Sigma(T)$ lies in the single coset
$d+\widetilde\Sigma(T)$, whose size is
$|\widetilde\Sigma(T)|=|\Sigma(T)|$.  The $|\Sigma(T)|$ fibres in this union are
nonempty and disjoint because they lie over distinct quotient directions.
Hence each is a singleton.  In particular $a\in\Sigma(T)$, so
$\overline a=0$, a contradiction.  Thus $\overline T$ has no nonzero singleton
directions.
\end{proof}

\end{document}
